% Detailed proofs for 04qed-ward-takahashi
\begin{derivation}[从\eqnrefs{eq:qed_inverse_propagator_gauge_covariance}到\eqnrefs{eq:qed_ward_takahashi_coordinate}]

正文\eqnrefs{eq:qed_inverse_propagator_gauge_covariance} 给出背景场变化时精确电子逆传播子的端点协变性。对任意无穷小规范函数 $\chi$ 比较等式两侧的一阶变化，便能把这种端点相位约束转写成顶角的坐标空间散度关系。

左端对背景场作泛函展开，
\begin{align*}
 S_{-\partial\chi}^{-1}(x,y)
 &=S_0^{-1}(x,y)
 -\int\dd^4 z\,
 \left.\frac{\delta S_A^{-1}(x,y)}{\delta A_\mu(z)}\right|_0
 \partial_\mu\chi(z)+O(\chi^2)\\
 &=S_0^{-1}(x,y)
 +\qe\int\dd^4 z\,\Gamma^\mu(x,y;z)\partial_\mu\chi(z)+O(\chi^2)\\
 &=S_0^{-1}(x,y)
 -\qe\int\dd^4 z\,[\partial_\mu^z\Gamma^\mu(x,y;z)]\chi(z)+O(\chi^2).
\end{align*}
右端的两个端点相位展开为
\begin{align*}
 U_\chi(x)S_0^{-1}(x,y)U_\chi^{-1}(y)
 &=S_0^{-1}(x,y)
 +\frac{\ii\qe}{\hbar}[\chi(x)-\chi(y)]S_0^{-1}(x,y)+O(\chi^2).
\end{align*}
两个端点值分别可写成
\begin{equation*}
 \chi(x)=\int\dd^4z\,\delta^{(4)}(z-x)\chi(z),
 \qquad
 \chi(y)=\int\dd^4z\,\delta^{(4)}(z-y)\chi(z).
\end{equation*}
所以右端一阶变化等于
\begin{equation*}
 \frac{\ii\qe}{\hbar}\int\dd^4z\,
 \bigl[\delta^{(4)}(z-x)-\delta^{(4)}(z-y)\bigr]
 S_0^{-1}(x,y)\chi(z).
\end{equation*}
与左端
\begin{equation*}
 -\qe\int\dd^4z\,
 [\partial_\mu^z\Gamma^\mu(x,y;z)]\chi(z)
\end{equation*}
相等。由于 $\chi(z)$ 是任意光滑测试函数，两个分布作用在任意测试函数上的积分都相等，只能有
\begin{equation*}
 \partial_\mu^z\Gamma^\mu(x,y;z)
 =-\frac{\ii}{\hbar}
 \bigl[\delta^{(4)}(z-x)-\delta^{(4)}(z-y)\bigr]
 S_0^{-1}(x,y),
\end{equation*}
即正文\eqnrefs{eq:qed_ward_takahashi_coordinate}。这里的两个 $\delta$ 接触项正是端点相位 $\chi(x)-\chi(y)$ 的分布表示：一个支撑在电子线的入射端点，另一个支撑在出射端点，二者符号相反。
\end{derivation}

\begin{derivation}[从\eqnrefs{eq:qed_ward_takahashi_coordinate}到\eqnrefs{eq:qed_ward_zero_momentum_vertex_derivative}]
按正文固定的 $\exp(-\ii p\cdot x/\hbar)$ 约定，把平移不变的逆传播子和三点顶角展开为
\begin{align*}
S^{-1}(x,y)
&=\int_p
 \ee^{-\ii p\cdot x/\hbar}
 \ee^{+\ii p\cdot y/\hbar}S^{-1}(p),\\
\Gamma^\mu(x,y;z)
&=\int_p\!\int_q
 \ee^{-\ii(p+q)\cdot x/\hbar}
 \ee^{+\ii p\cdot y/\hbar}
 \ee^{+\ii q\cdot z/\hbar}
 \Gamma^\mu(p+q,p),
\end{align*}
其中 $\int_p\equiv\int \dd^4 p/(2\pi\hbar)^4$。正文\eqnrefs{eq:qed_ward_takahashi_coordinate} 左边的 $z$ 导数只作用在最后一个 Fourier 相位上，因此
\begin{align*}
\partial_\mu^z\Gamma^\mu(x,y;z)
&=\int_p\!\int_q
 \ee^{-\ii(p+q)\cdot x/\hbar}
 \ee^{+\ii p\cdot y/\hbar}
 \ee^{+\ii q\cdot z/\hbar}
 \frac{\ii q_\mu}{\hbar}
 \Gamma^\mu(p+q,p).
\end{align*}
两个接触项分别使用
\begin{equation*}
\delta^{(4)}(z-x)=\int_q\ee^{+\ii q\cdot(z-x)/\hbar},
\qquad
\delta^{(4)}(z-y)=\int_q\ee^{+\ii q\cdot(z-y)/\hbar}.
\end{equation*}
第一项给
\begin{align*}
\delta^{(4)}(z-x)S^{-1}(x,y)
=\int_p\!\int_q
 \ee^{-\ii(p+q)\cdot x/\hbar}
 \ee^{+\ii p\cdot y/\hbar}
 \ee^{+\ii q\cdot z/\hbar}S^{-1}(p),
\end{align*}
而第二项把逆传播子的积分变量改写成 $p+q$ 后成为
\begin{align*}
\delta^{(4)}(z-y)S^{-1}(x,y)
=\int_p\!\int_q
 \ee^{-\ii(p+q)\cdot x/\hbar}
 \ee^{+\ii p\cdot y/\hbar}
 \ee^{+\ii q\cdot z/\hbar}S^{-1}(p+q).
\end{align*}
于是坐标空间恒等式两侧具有同一个 Fourier 核。比较系数，
\begin{equation*}
\frac{\ii}{\hbar}q_\mu\Gamma^\mu(p+q,p)
=-\frac{\ii}{\hbar}\left[S^{-1}(p)-S^{-1}(p+q)\right],
\end{equation*}
约去共同的 $\ii/\hbar$，得到正文\eqnrefs{eq:qed_ward_takahashi_momentum}：
\begin{equation*}
 q_\mu\Gamma^\mu(p+q,p)
 =S^{-1}(p+q)-S^{-1}(p).
\end{equation*}

再取 $q\to0$，并使用正文已经说明的可微条件。右边作 Taylor 展开：
\begin{equation*}
 S^{-1}(p+q)-S^{-1}(p)
 =q_\mu\frac{\partial S^{-1}(p)}{\partial p_\mu}+O(q^2).
\end{equation*}
左边的精确顶角在同一点展开为
\begin{equation*}
 q_\mu\Gamma^\mu(p+q,p)
 =q_\mu\Gamma^\mu(p,p)+O(q^2).
\end{equation*}
两边对任意足够小的四矢量 $q_\mu$ 都相等，因此比较每个独立的一阶系数得到
\begin{equation*}
 \Gamma^\mu(p,p)=\frac{\partial S^{-1}(p)}{\partial p_\mu},
\end{equation*}
即正文\eqnrefs{eq:qed_ward_zero_momentum_vertex_derivative}。
\end{derivation}

\begin{derivation}[从\eqnrefs{eq:qed_ward_zero_momentum_vertex_derivative}到\eqnrefs{eq:qed_ward_z1_equals_z2}]

完整 1PI 顶角与完整电子逆传播子满足
\[
q_\mu\Gamma_B^\mu(p+q,p)
=
S_B^{-1}(p+q)-S_B^{-1}(p).
\]
令 \(q\to0\)，对 \(q_\mu\) 作一阶展开即得正文
\eqnrefs{eq:qed_ward_zero_momentum_vertex_derivative}：
\[
\Gamma_B^\mu(p,p)
=
\frac{\partial S_B^{-1}(p)}{\partial p_\mu}.
\]
这个关系作用在完整裸量上，因此任何紫外发散都必须在等式两边以相同方式出现。

要从这个恒等式抽取重整化常数，最干净的做法是只投影它的局域紫外发散部分，而不把方案依赖的有限项混进来。由表面发散度和 Lorentz 协变性，电子二点 1PI 核的局域紫外发散最多具有
\begin{equation*}
 \bigl[S_B^{-1}(p)\bigr]_{\rm UV,loc}
 =A_{\rm UV}\,\slashed p
 +B_{\rm UV}\,m_Rc,
\end{equation*}
其中 $A_{\rm UV},B_{\rm UV}$ 可以含 $1/\epsilon$ 极点，但不再依赖外动量；更高次的 $p$ 幂会对应不可重整化的局域算符。对 $p_\mu$ 求导后，质量型发散自动消失：
\begin{equation*}
 \frac{\partial}{\partial p_\mu}
 \bigl[S_B^{-1}(p)\bigr]_{\rm UV,loc}
 =A_{\rm UV}\gamma^\mu.
\end{equation*}

另一方面，零光子动量处的三点 1PI 顶角在 QED 中只有同维数的局域 Dirac 结构
\begin{equation*}
 \bigl[\Gamma_B^\mu(p,p)\bigr]_{\rm UV,loc}
 =C_{\rm UV}\gamma^\mu.
\end{equation*}
把这两个局域发散部分代入精确零动量 Ward--Takahashi 恒等式
\begin{equation*}
 \Gamma_B^\mu(p,p)
 =\frac{\partial S_B^{-1}(p)}{\partial p_\mu},
\end{equation*}
于是恒等式在唯一允许的局域基矢 $\gamma^\mu$ 上给出：
\begin{equation*}
 C_{\rm UV}\gamma^\mu=A_{\rm UV}\gamma^\mu
 \quad\Longrightarrow\quad
 C_{\rm UV}=A_{\rm UV}.
\end{equation*}

现在把反项写成
\begin{equation*}
 \delta_2\,\bar\psi_R\,\ii\hbar c\slashed\partial\psi_R,
 \qquad
 -\delta_1 q_Rc\,\bar\psi_R\gamma^\mu\psi_R A_{\mu,R},
 \qquad
 \delta_i\equiv Z_i-1.
\end{equation*}
重整化条件要求反项逐个抵消上述局域紫外系数，因此在同一符号约定下
\begin{equation*}
 \delta_2=-A_{\rm UV},
 \qquad
 \delta_1=-C_{\rm UV}.
\end{equation*}
由 $C_{\rm UV}=A_{\rm UV}$ 立即有
\begin{equation*}
 \delta_1=\delta_2,
 \qquad
 \boxed{Z_1=Z_2.}
\end{equation*}
质量反项不参与这一投影，因为 $\partial(m_Rc)/\partial p_\mu=0$；有限顶角与有限自能导数则继续由同一个 Ward--Takahashi 恒等式彼此约束。

因此，在保持局域 $U(1)$ Ward 恒等式的调节与重整化方案中，$Z_1=Z_2$ 是全阶约束。结合裸电荷与重整化电荷的关系，独立的电荷重整化信息只剩光子场因子 $Z_3$。
\end{derivation}

\begin{derivation}[从\eqnrefs{eq:qed-vacpol-quadratic-kernel-dp68}到\eqnrefs{eq:qed_vacuum_polarization_transversality_from_1pi}]

正文\eqnrefs{eq:qed-vacpol-quadratic-kernel-dp68} 把真空极化定义成光子有效作用量中的量子二次核。对它施加无穷小规范变换后，任意规范函数 \(\chi(q)\) 前的系数必须消失；这会直接限制 \(\Pi^{\mu\nu}\) 的纵向分量。

Fourier 空间的无穷小变换为
\begin{equation*}
\delta A_\mu(q)=-\frac{\ii}{\hbar}q_\mu\chi(q).
\end{equation*}
对正文\eqnrefs{eq:qed-vacpol-quadratic-kernel-dp68} 作一阶变分。二次核满足 $\Pi^{\mu\nu}(q)=\Pi^{\nu\mu}(-q)$；第二项作 $q\mapsto-q$ 后可与第一项合并，严格得到
\begin{equation*}
\delta\Gamma_{\Pi}^{(2)}
=-\frac{\ii}{\hbar}
\int\frac{\dd^4 q}{(2\pi\hbar)^4}\,
\chi(-q)q_\mu\Pi^{\mu\nu}(q)A_\nu(q).
\end{equation*}
规范函数 \(\chi\) 与测试场 \(A_\nu\) 可以独立选择，而量子有效作用量必须对任意 \(\chi\) 保持不变，因此只有
\begin{equation*}
q_\mu\Pi^{\mu\nu}(q)=0
\end{equation*}
才能使变分恒为零，这就是正文\eqnrefs{eq:qed_vacuum_polarization_transversality_from_1pi}。
\end{derivation}
